%------------------------------------------------------------------------------%
\documentclass[a4paper,12pt,fleqn]{article}

\usepackage{calculo_varias_variaveis-1}

\ShowAnswers

%------------------------------------------------------------------------------%
\begin{document}

\makeHeader{CFVVI}{Prova 2 -- Turma 6}{28/09/2026}

\bigskip
\textbf{Atenção}: a aplicação das regras do produto ou do quociente em casos
nos quais um dos termos é constante \textbf{será considerada incorreta},
por evidenciar compreensão inadequada das regras de derivação.

% 01
%------------------------------------------------------------------------------%
\exercise{25\%}
Calcule a derivada parcial mista
\(
  f_{xy}(x, y)
\)
da função
\(
  f(x, y) = ye^{xy} - x\sin\left(x^2 + y\right),
\)
na ordem indicada pela notação

\begin{answer}
  \begin{align*}
    f_x(x, y)
    & = \D{f}{x} \\
    & = \D{}{x}\left[ye^{xy} - x\sin\left(x^2 + y\right)\right] \\
    & = \D{}{x}\left[ye^{xy}\right]
      - \D{}{x}\left[x\sin\left(x^2 + y\right)\right] \\
    & = y\D{}{x}\left[e^{xy}\right] -
        \left(
          \D{x}{x}\sin\left(x^2 + y\right)
        + x\D{}{x}\left[\sin\left(x^2 + y\right)\right]
        \right)
    \\
    & = ye^{xy}\D{}{x}(xy)
      - \sin\left(x^2 + y\right)
      - x\cos\left(x^2 + y\right)\D{}{x}\left(x^2 + y\right)
    \\
    & = ye^{xy}y
      - \sin\left(x^2 + y\right)
      - x\cos\left(x^2 + y\right)2x
    \\
    & = y^2e^{xy}
      - \sin\left(x^2 + y\right)
      - 2x^2\cos\left(x^2 + y\right)
  \end{align*}
  \begin{align*}
    f_{xy}(x, y)
    & = \DM{f}{y}{x} \\
    & = \D{}{y}\left(\D{f}{x}\right)
    \\
    & = \D{}{y}\left[
          y^2e^{xy}
          - \sin\left(x^2 + y\right)
          - 2x^2\cos\left(x^2 + y\right)
        \right]
    \\
    & = \D{}{y}\left[ y^2e^{xy} \right]
      - \D{}{y}\left[ \sin\left(x^2 + y\right) \right]
      - \D{}{y}\left[ 2x^2\cos\left(x^2 + y\right) \right]
    \\
    & = \D{y^2}{y}e^{xy} + y^2\D{e^{xy}}{y}
      - \D{}{y}\left[ \sin\left(x^2 + y\right) \right]
      - 2x^2\D{}{y}\left[ \cos\left(x^2 + y\right) \right]
    \\
    & = 2ye^{xy} + y^2e^{xy}\D{xy}{y}
      - \cos\left(x^2 + y\right) \D{}{y}\left[ x^2 + y \right]
      + 2x^2\sin\left(x^2 + y\right)\D{}{y}\left[ x^2 + y \right]
    \\
    & = 2yxe^{xy} + y^2e^{xy}
      - \cos\left(x^2 + y\right)
      + 2x^2\sin\left(x^2 + y\right)
  \end{align*}
\end{answer}

% 02
%------------------------------------------------------------------------------%
\exercise{25\%}
  Determine se o limite
  \(
    \lim_{(x, y)\to(0, 0)} \frac{x^2-y^2}{x^2+y^2}
  \)
  existe

\begin{answer}
  \[
    f(x, y) = \frac{x^2-y^2}{x^2+y^2}
  \]
  Caminho $y=0$
  \[
    f(x, 0)
    = \frac{x^2-y^2}{x^2+y^2} \evalat{y=0}{}
    = \frac{x^2}{x^2}
    = 1
    \qquad\forall x\neq 0
  \]
  \[
    L_1
    = \lim_{(x, y)\to(0, 0)} \frac{x^2-y^2}{x^2+y^2} \evalat{y=0}{}
    = \lim_{x \to 0} \frac{x^2-0^2}{x^2+0^2}
    = \lim_{x \to 0} 1
    = 1
  \]
  Caminho $x = 0$
  \[
    f(0, y)
    = \frac{x^2-y^2}{x^2+y^2} \evalat{x=0}{}
    = \frac{-y^2}{y^2} 4
    = -1
    \qquad\forall y\neq 0
  \]
  \[
    L_2
    = \lim_{(x, y)\to(0, 0)} \frac{x^2-y^2}{x^2+y^2} \evalat{x=0}{}
    = \lim_{y \to 0} \frac{0^2-y^2}{0^2+y^2}
    = \lim_{y \to 0} -1
    = -1
  \]
  Os dois caminhos conduzem ao ponto $(0, 0)$,
  mas produzem valores diferentes então o limite não existe.
\end{answer}

% 03
%------------------------------------------------------------------------------%
\exercise{25\%}
Dada a equação
\(
  x^2 + xy + y^2 = 7,
\)
determine \(\frac{dy}{dx}\) por derivação implícita e calcule seu valor no ponto $(2, 3)$

\begin{answer}
  \textbf{Método 1}
  Pela fórmula temos
  \[
    \frac{dy}{dx} = - \frac{F_x}{F_y}
  \]
  com
  \[
    F(x, y) = x^2 + xy + y^2
  \]
  \[
    F_x(x, y)
    = \D{F}{x}
    = \D{}{x}\left(x^2 + xy + y^2\right)
    = 2x + y
  \]
  \[
    F_y(x, y)
    = \D{F}{y}
    = \D{}{y}\left(x^2 + xy + y^2\right)
    = x + 2y
  \]
  \[
    \frac{dy}{dx}
    = -\frac{F_x}{F_y}
    = -\frac{2x + y}{x + 2y}
  \]

  \textbf{Método 2}

  Como $y$ é uma função de $x$,
  \(
    y = y(x)
  \)
  \begin{align*}
    x^2 + xy + y^2                                                & = 7                  \\
    x^2 + xy(x) + \left(y(x)\right)^2                             & = 7                  \\
    \frac{d}{dx} \left[ x^2 + xy(x) + \left(y(x)\right)^2 \right] & = \frac{d7}{dx}      \\
    2x + \frac{dx}{dx}y + x\frac{dy}{dx} + 2y\frac{dy}{dx}        & = 0                  \\
    2x + y + x\frac{dy}{dx} + 2y\frac{dy}{dx}                     & = 0                  \\
    \frac{dy}{dx}\left(x + 2y\right) + 2x + y                     & = 0                  \\
    \frac{dy}{dx}\left(x + 2y\right)                              & = - 2x - y           \\
    \frac{dy}{dx}                                                 & = \frac{-2x-y}{x+2y}
  \end{align*}

  \textbf{No ponto $(1,2)$}

  \[
    \frac{dy}{dx}(2, 3)
    = -\frac{2x + y}{x + 2y} \evalat{(2, 3)}{}
    = -\frac{2\times 2 + 3}{2 + 2 \times 3}
    = -\frac{7}{8}
  \]
\end{answer}

% 04
%------------------------------------------------------------------------------%
\exercise{25\%}
Dados
$
  z=u^2+e^v
$,
$
  u=x^2+y
$ e
$
  v=xy
$.
Utilize a regra da cadeia para calcular
\(
  \D{z}{x}(1, 2)
\)
\begin{answer}
\[
  \D{z}{x} = \D{z}{u} \D{u}{x} + \D{z}{v} \D{v}{x}
\]
\[
  \D{z}{u} = 2u
\]
\[
  \D{z}{v} = e^v
\]
\[
  \D{u}{x} = 2x
\]
\[
  \D{v}{x} = y
\]
\begin{align*}
  \D{z}{x}
  & = \D{z}{u} \D{u}{x} + \D{z}{v} \D{v}{x} \\
  & = 2u(2x) + e^v y                  \\
  & = 4x\left(x^2 + y\right)+ye^{xy}
\end{align*}
\begin{align*}
  \D{z}{x}(1, 2)
  & = \left(4x\left(x^2 + y\right)+ye^{xy}\right) \evalat{(1, 2)}{} \\
  & = 4(1)(1^2+2)+2e^{1 \times 2} \\
  & = 12+2e^2
\end{align*}
\end{answer}

%------------------------------------------------------------------------------%
\end{document}
%------------------------------------------------------------------------------%
